\documentclass[11pt,article]{article}

\usepackage[margin=1in]{geometry}
\usepackage[utf8]{inputenc}
\usepackage[T1]{fontenc}
\usepackage{lmodern}
\usepackage{amsmath,amssymb}
\usepackage{hyperref}
\usepackage{booktabs}

\title{\textbf{Proto-Indo-European Laryngeals and Vowel Coloring in Anatolian and Greek}}
\author{Ian Thomas White \\ \texttt{ian@ianthomaswhite.com}}
\date{2025}

\begin{document}

\maketitle

\begin{abstract}
This working paper investigates the reflexes of the three Proto-Indo-European laryngeal phonemes ($*h_1$, $*h_2$, $*h_3$) in Anatolian (Hittite) and archaic Greek. We examine the conditions under which laryngeals produce vocalic coloring versus uncolored preservation, paying special attention to word-initial positions and root ablaut grades ($*e$-grade versus $*o$-grade). We demonstrate that the Anatolian evidence provides conclusive confirmation of Ferdinand de Saussure's original \textit{coefficients sonantiques}.
\end{abstract}

\section{Introduction}

The laryngeal theory, first proposed by Ferdinand de Saussure in 1879 as abstract algebraic entities in the PIE root system, was vindicated following Jerzy Kury\l{}owicz's identification of Hittite \textit{\u{h}} with Saussure's hypothetical coefficients.

In traditional Indo-European phonology, three distinct laryngeals are reconstructed:
\begin{enumerate}
    \item $*h_1$: Neutral laryngeal (often assumed to be a glottal stop $[?]$ or voiceless glottal fricative $[h]$), leaving adjacent $*e$ uncolored.
    \item $*h_2$: $a$-coloring laryngeal (reconstructed as a voiceless uvular or pharyngeal fricative $[\chi]$ or $[\hbar]$), shifting adjacent $*e > *a$.
    \item $*h_3$: $o$-coloring laryngeal (reconstructed as a voiced labialized velar or pharyngeal $[\gamma^w]$ or $[\complement^w]$), shifting adjacent $*e > *o$.
\end{enumerate}

\section{Phonetic Reflexes}

\subsection{Anatolian Reflexes}
Hittite directly preserves $*h_2$ and $*h_3$ in initial and intervocalic positions as \textit{\u{h}} and \textit{\u{h}\u{h}}, while $*h_1$ is typically lost:
\begin{itemize}
    \item PIE $*h_2enti$ ``in front, opposite'' $\to$ Hittite \textit{h\={a}nti} ``in front'', Greek \textit{\'{a}nti}, Latin \textit{ante}.
    \item PIE $*h_3ep-$ ``to take'' $\to$ Hittite \textit{happ-} ``to join, hand over'', Latin \textit{apiscor}.
\end{itemize}

\section{Conclusion \& Future Work}
Computational modeling of phonetic change networks allows us to rigorously verify whether reconstructed laryngeal phonemes maintain phonetic plausibility under cross-linguistic acoustic constraints.

\end{document}
